\clearpage
\section{系统设计}

\subsection{总体架构}
系统由服务层、记录与验证层、观测层组成。客户端经 HTTP 访问键值集群或事务服务。五个 Raft 节点通过节点间 RPC 复制命令，各自使用独立的 SQLite 数据卷。事务服务在独立进程中执行多键事务，由同一持久卷保存 WAL 和 SQLite 版本表。两套服务的操作记录采用共享数据契约，分别进入一致性检查。图~\ref{fig:architecture} 给出组件关系。

\begin{figure}[htbp]
\centering
\begin{tabular}{@{}c@{\hspace{13mm}}c@{}}
\ReportDiagramNode{.35\linewidth}{键值客户端} &
\ReportDiagramNode{.35\linewidth}{事务客户端} \\
\small ↓ HTTP 键值接口 & \small ↓ HTTP 事务接口 \\
\ReportDiagramNode{.35\linewidth}{Raft 键值集群\\五节点通过 RPC 复制命令} &
\ReportDiagramNode{.35\linewidth}{独立事务服务\\RC 与 SI 多版本读写} \\
\small ↓ 各节点独立持久化 & \small ↓ 服务本地持久化 \\
\ReportDiagramNode{.35\linewidth}{每节点 SQLite\\任期、日志和快照} &
\ReportDiagramNode{.35\linewidth}{事务卷\\WAL 与 SQLite 版本表} \\
\noalign{\vskip 5pt}
\multicolumn{2}{c}{\small 调用历史和版本记录\ReportDiagramArrow 离线验证\ReportDiagramArrow 判定归档} \\
\multicolumn{2}{c}{\small 状态、容器事件与判定归档\ReportDiagramArrow Observer\ReportDiagramArrow 浏览器}
\end{tabular}

\caption{系统组成与主要数据流}
\label{fig:architecture}
\end{figure}

Observer 并行采样 Raft 节点，读取容器日志与运行归档，并通过本机端口查询事务服务。浏览器经同源接口访问 Observer。运维操作通过运行器执行，项目锁协调同一 Compose 集群上的互斥操作；协议核心继续由节点进程驱动。各模块的状态归属见表~\ref{tab:ownership}。

\begin{table}[htbp]
\centering
\caption{主要模块及其职责}
\label{tab:ownership}
\begin{tabularx}{\linewidth}{@{}l p{5cm} X@{}}
\toprule
模块 & 源码位置 & 职责与状态 \\
\midrule
Raft 协议与适配 & \path{src/raft/} & 选举、日志复制、提交和读屏障；经 SQLite 和 HTTP 适配器执行持久化与传输。 \\
事务执行与恢复 & \path{src/transaction/} & 事务上下文、私有写集、锁、版本链和 WAL；启动时恢复已提交状态。 \\
记录契约与验证 & \path{src/contracts/}\newline\path{src/verification/} & 定义历史记录；从调用及版本证据计算顺序、依赖和业务判定。 \\
观测与交互 & \path{src/observer/}、\path{web/} & 汇集状态、容器事件、事务记录和判定，提供页面查询与回放。 \\
\bottomrule
\end{tabularx}
\end{table}

\par\vbox{%
\subsection{节点角色与数据结构}
Follower 接收 Leader 的追加日志请求，保存任期、投票和复制进度所需的本地日志；选举计时到期后，具备投票资格的节点转为 Candidate。Candidate 自投一票并向其他投票节点请求投票。取得当前配置的多数票后，节点成为 Leader，管理客户端请求及各副本的复制进度。\par
}

\begin{table}[htbp]
\centering
\caption{关键结构与顺序标识}
\label{tab:structures}
\begin{tabularx}{\linewidth}{@{}l >{\raggedright\arraybackslash}X >{\raggedright\arraybackslash}X@{}}
\toprule
结构 & 主要字段 & 用途 \\
\midrule
任期与投票 & \texttt{term}、\texttt{votedFor} & 任期标识领导周期；任期和投票结果持久化后才对外应答。 \\
日志和应用水位 & 条目任期、命令、绝对索引；\texttt{commitIndex}、\texttt{lastApplied} & 提交水位决定可应用前缀；应用水位标记状态机已经执行的位置。 \\
复制进度与快照 & \texttt{nextIndex}、\texttt{matchIndex}；快照边界索引、任期、状态和配置 & 定位待复制后缀；落后节点先安装快照，再按绝对索引追赶。 \\
事务及版本链 & 事务 ID、隔离级别、快照序号、写集；版本 ID、\texttt{commitSeq}、值或墓碑 & 将事务私有修改与已提交版本分开；提交序号决定逻辑可见性。 \\
WAL 记录 & \texttt{lsn}、\texttt{prevLsn}、类型及负载 & LSN 定位物理动作，前驱指针连接同一事务的记录。 \\
操作历史 & 操作 ID、调用及返回时刻、结果、读写版本 & 保存客户端视角的顺序约束以及事务依赖所需的版本身份。 \\
\bottomrule
\end{tabularx}
\end{table}

Raft 的日志索引表示一条复制命令在集群日志中的位置；事务的提交序号描述已提交版本的先后；WAL 的 LSN 则标识物理动作。事务的多条 WAL 记录可以对应一次提交，多个键的最终版本在同一提交序号下发布。

\subsection{Raft 选举与读写流程}
\subsubsection{选举与日志追赶}
选举超时后，Candidate 增加任期并持久化自投票。接收端处理请求任期，先比较最后一条日志的任期，再比较末尾索引，并检查本任期的投票状态。投票结果持久化后应答；取得配置多数票后，新 Leader 追加当前任期空操作。

追加日志时，Follower 对照请求的前序索引与任期。匹配后持久化新增或替换的后缀；提交水位以 Leader 公告与本次匹配末尾的较小值为限。任期冲突时，Follower 返回冲突任期在本地的首个索引。Leader 用仍与当前探针匹配的失败应答修正发送位置，并保留已确认的复制进度。发送位置落入已压缩前缀时改发快照。

\subsubsection{写入与读取}
写命令在 Leader 按批追加到本地日志，再发送给 Follower。Leader 根据各节点已匹配的位置寻找达到多数且属于当前任期的条目。提交推进后，状态机按索引应用新增前缀，客户端随后收到结果。图~\ref{fig:raft-paths} 对照了选举、写入和读屏障。

\begin{figure}[htbp]
\centering
\ReportDiagramNode{.26\linewidth}{选举超时\\持久化任期和自投票}\ReportDiagramArrow
\ReportDiagramNode{.26\linewidth}{RequestVote\\比较末尾日志并投票}\ReportDiagramArrow
\ReportDiagramNode{.26\linewidth}{多数票选出 Leader\\追加本任期空操作}\\[5pt]
\ReportDiagramNode{.26\linewidth}{写入请求\\本地追加并持久化}\ReportDiagramArrow
\ReportDiagramNode{.26\linewidth}{AppendEntries\\对齐前缀并收集确认}\ReportDiagramArrow
\ReportDiagramNode{.26\linewidth}{本任期条目获多数\\推进提交、应用并应答}\\[5pt]
\ReportDiagramNode{.26\linewidth}{读取请求\\登记读编号和读索引}\ReportDiagramArrow
\ReportDiagramNode{.26\linewidth}{本任期空操作已应用\\多数确认该读编号}\ReportDiagramArrow
\ReportDiagramNode{.26\linewidth}{本地应用达到读索引\\返回键值或读屏障}

\caption{Raft 选举、写入和读屏障的关键路径}
\label{fig:raft-paths}
\end{figure}

Leader 将读编号随追加日志 RPC 发送，成功回复才累计读确认。读屏障须满足本任期空操作已应用、读编号获多数确认、本地应用达到发起时的提交索引。返回的读索引取该提交索引与本任期空操作位置的较大值。

\par\vbox{%
默认读由 Leader 在屏障后执行。安全 Follower 读向 Leader 请求新鲜屏障，再等待本地应用到返回的索引；超时或拿不到屏障时返回 \path{NO_QUORUM}。显式本地读直接返回已应用值和应用索引，分区期间可读到旧值。\par
}

\subsection{事务提交与恢复流程}
事务开始时分配 ID，快照隔离事务同时记录开始时的提交序号。读取先检查本事务私有写集，然后按隔离级别选择已提交版本。每次写入记录一条 \texttt{UPDATE} 并在 SQLite 中物化为待提交版本。提交时引擎完成快照隔离的写冲突检查，持久化 \texttt{COMMIT} 后发布每个键的最终版本，最后以 \texttt{END} 结束该事务的日志链。

\begin{figure}[htbp]
\centering
\ReportDiagramNode{.26\linewidth}{提交请求\\完成冲突检查并分配序号}\ReportDiagramArrow
\ReportDiagramNode{.26\linewidth}{同步 COMMIT\\保存提交决定}\ReportDiagramArrow
\ReportDiagramNode{.26\linewidth}{SQLite 发布版本\\追加 END 并返回}\\[5pt]
\ReportDiagramNode{.26\linewidth}{分析\\按检查点恢复事务链}\ReportDiagramArrow
\ReportDiagramNode{.26\linewidth}{重做\\应用游标跳过已物化动作}\ReportDiagramArrow
\ReportDiagramNode{.26\linewidth}{撤销与收尾\\记录 CLR 或补齐 END}

\caption{事务提交与启动恢复}
\label{fig:recovery}
\end{figure}

最近完整检查点保存活动事务的日志链尾和当时的应用游标。分析从该检查点之后重建尚未结束的事务；重做从检查点记录的游标之后扫描，SQLite 用当前游标跳过已物化动作。已持久化 \texttt{COMMIT} 的事务恢复版本并补齐收尾记录；其余事务沿前驱链撤销，补偿记录（CLR）保存继续撤销的位置。图~\ref{fig:recovery} 展示两条路径如何共用 WAL。
